% !TeX root = ../papers/nc-p3.tex
\section{二つの反対隅の四行列恒等式からの全正則窓}
\label{sec:opposite-four-matrix-corners}
\begin{theorem}\label{thm:opposite-four-matrix-corners}
有限四辺の標準正則窓で一隅が四行列の恒等式を満たすとする。
指定隅で次元3の単純は全位置で次元3かつ頂点と両向きに直交する。
次元2の場合に残り得る非均衡型は $(2,2,2,3),(2,3,2,3),(2,3,3,3)$ の巡回移動だけである。
任意の二隅に同じ恒等式があれば全単純は均衡する。
さらにその二隅が反対位置なら全窓は中心有限な合成整域で、定理\ref{thm:four-matrix-all-regular}の
32次までの判定がこの枝の全正則点を認識する。
\end{theorem}

\begin{proof}[次元型の有限絞り込み]
指定隅の集合を $J$ とする。薄い場合と、指定隅で次元4に達する場合は既出なので、
ある指定隅で次元2または3の非薄単純 $S$ を扱う。$T$ をその一般の異なる次数倍率とし、両者間の第一Extは両向きに零である。
一指定隅の次元 $d$ が $2d>4$ なら行列単位の議論で直接従い、$d=2$ でも $3d>4$ のため
定理\ref{thm:one-sided-balanced-small-simples}冒頭で一般の中間倍率をいくつでも挿入する議論が使える。
この段階には均衡性を使わない。$E=\sum d_i h_i$ は二端間の第二Ext次元である。
以下 $2d_t>4$ の指定隅 $t$ がある場合を「二端だけで超過」と書く。
各頂点での $a_i,b_i,h_i$ の上界・下界は前節と同じである。$E\ge\sum a_t b_t$ は
二端だけで超過するなら全四位置の和、そうでなければ $t\in J$ の和について成立する。

$i,j=i+1,k=i+2$ を取り、$A=a_i,B=b_k$ とする。この段落の $B$ は重複度である。
$S$ を下から $k_i^A$ で、$T$ を上から $k_k^B$ で普遍拡張し $U,V$ を作る。
$K=h_j+a_j+b_j$ とすると
\[
 u=a_j+\max(0,4A-K),\qquad w=b_j+\max(0,4B-K)
\]
は $\operatorname{Ext}^1(U,k_j),\operatorname{Ext}^1(k_j,V)$ の次元下界である。
$k_i$ 部分に制限して検出される前者の次元は少なくとも $u-a_j$、$k_k$ 商へ移して検出される後者は
少なくとも $w-b_j$ である。いずれも普遍拡張の完全列による。
二端だけで超過するか $j\in J$ なら $u'=u,w'=w$ とする。それ以外では
$i\in J$ なら $u'=u-a_j$、そうでなければ0とし、$k\in J$ なら $w'=w-b_j$、そうでなければ0とする。
$L_i$ は二端だけで超過するか $i\in J$ なら0、それ以外なら $A b_i$ とし、
$L_k$ も同じ規約で $0$ または $B a_k$ とする。次の不等式を使う。
\begin{align*}
 u'b_j&\le E+A(h_i+a_i)+L_i,& a_jw'&\le E+B(h_k+b_k)+L_k,\\
 uw-(u-u')(w-w')&\le E+A(h_i+a_i)+B(h_k+b_k)+6AB+L_i+L_k-\max(A r_a,B r_b).
\end{align*}
$r_a,r_b$ は後述の矢の二段積の有効次元の下界である。

左辺を説明する。中間 $k_j$ のYoneda積が消えるテンソルを階数分解し、定理\ref{thm:four-matrix-balanced-simples}と
同様に三段表現に組む。中間の各重複方向には頭 $S$ から到達でき、どの非零方向からもsocle $T$ に到達できる。
二端だけで超過するか $j\in J$ なら従来通り核は零である。
$i\in J$ の場合、$k_i^A$ に制限したテンソルが非零なら中間の少なくとも一つの方向に
$k_i$ ブロックを通って到達する。$U$ は頭 $S$ で生成されるから、そこから $S,k_i,T$ を通る非零積が得られ、
指定隅の三次元数の和は少なくとも $2+1+2$ となる。$k\in J$ なら商側で同じである。
従って核は両側の指定された制限核のテンソル積に入る。必要なら実際の空間から
上の下界個の独立な類を、制限の所定の階数を保って選べる。これを $(U,T),(S,V),(U,V)$ の順で使うと各左辺を得る。

上界では $\dim\operatorname{Ext}^1(U,T)\le L_i$ を使う。普遍拡張によりこの群は
$k_i$ を通る $A b_i$ 次元のYoneda写像の核であり、指定隅によっては単射となるからである。
同様に $\dim\operatorname{Ext}^1(S,V)\le L_k$。完全列により最初の二つの右辺を得る。
また $\dim\operatorname{Ext}^1(U,V)\le L_i+L_k$ である。
実際 $\operatorname{Ext}^1(U,k_k)=\operatorname{Ext}^1(S,k_k)$（$i,k$ は反対位置）で、
$S,V$ の完全列で $k_k$ を通る積が単射でも $U$ への制限で核が生じ得る。
その追加核の次元は $\operatorname{Ext}^2(S,T)\to\operatorname{Ext}^2(U,T)$ の核との交叉である。
両積の像の交叉は、二端だけで超過するか $i$ または $k$ が指定隅なら零である。
実際二中間型の和での関係はその指定位置を通る非零積を与える。どちらも指定されない場合は単なる
$\dim\operatorname{Ext}^1(U,T)+B a_k$ の上界を使えばよい。
$\chi(U,V)=E+A h_i+B h_k+6AB$ に、逆向き第一Extの上界を加え、CY双対を使えば最後の右辺を得る。

その逆向きの改善分を示す。$f=d_{i+3}$ と置く。二段 $k_k\to S\to k_i$ のYoneda積は、
$f$ 次元の中間空間を通る列解・行解を掛け、六関係で評価する写像である。
実際三段の上三角矢を掛ければ関係の非対角隅にこの積が出る。
頂点間の第二Extは六関係の双対であり、元の $d_i$ 行・$d_k$ 列では消える。
全行空間の次元を $p=d_i+A$、全列空間の次元を $l=d_k+B$ とする。
列の全てとの積が消える行の実効余次元を $r_a^{\rm act}$、行の全てとの積が消える列の余次元を
$r_b^{\rm act}$ とする。これらは倍率を変えても同じである。
$\operatorname{Ext}^1(V,k_i)$ は、$k_k$ との第一Extの消滅により、この行の積の核に入り
次元は $a_i-r_a^{\rm act}$ 以下である。反対端も同様。
従って $\dim\operatorname{Ext}^1(V,U)\le A a_i+B b_k-\max(A r_a^{\rm act},B r_b^{\rm act})$ となる。

$l>4+\max(0,3(f-3))$ なら $r_a=\max(0,p-2(f-2))$ を、他は $r_a=0$ を下界に使える。
実際全列を消す階数2の行があれば、その行と全列の積の空間は、四文字対二文字の行列のうち六関係を消す部分に入る。
その空間の次元は高々4である。5以上なら階数1のSegreとの交叉が正次元となり有限辺に反する。
この行との行列積そのものが零の列は幅 $f-2$ の行列で、列解が階数1を含まないため次元は高々
$\max(0,3(f-3))$ である。従って上の $l$ なら全列の関係評価を消す行は全て階数3以上である。
四対 $f$ 行列の階数2以下の余次元 $2(f-2)$ でその空間の次元を抑えればよい。
列側は $p,l$ を交換して同じ $r_b$ の下界を使う。

\path{computations/verify_opposite_four_matrix_corners.py} はこれらの整数不等式を全範囲で検算する。
指定隅を0、$d_0=2,3$ とし、向きの反転まで $2\le x=d_1\le z=d_3\le3(d_0-1)$、
$2\le d_2\le3(x-1)$ だけを列挙すればよい。
$d_0=3$ では $(3,3,3,3)$ のみ、$d_0=2$ では均衡型と定理中の三混合型の向きだけが残る。
混合三型の全ての指定二位置で同じ不等式は矛盾する。均衡三次元では、指定隅の二端だけで超過するので
前節の最短区間の議論が使え、$a_i,b_i\le3$ の下で少なくとも片側が全て零となる。
補題\ref{lem:periodic-tail-vertex-bridges}の薄い単純との三段で他方も消える。
\end{proof}

\begin{proof}[反対二位置での全次数判定]
二次元均衡単純を取る。$a_i,b_i\le1$、$E=0$ で、指定二位置では $a_i b_i=0$ である。
両向きが存在し、$a,b$ の非零位置の和集合が二位置以上なら、巡回して相異なる連続した非零位置
$i,j$ を $a_i>0,b_j>0$ となるよう選べる。実際そのような移行がなければ、ある $a$ の次から
全非零位置が $a$ のみとなり、出発点に戻っても $b$ を持てない。
その間の開区間は全て両方零で、閉区間は指定位置を含む（反対二位置の補集合に隣接二頂点はない）。
定理\ref{thm:one-sided-balanced-small-simples}の最短区間の橋でこの指定位置を通る非零積を作れる。
この論証は両端だけでなく途中の指定位置にも使える。非零な長さ $l$ の経路の各中間で
頂点冪等元を挿入し、半単純商のその頂点単純のブロックを選ぶ。中間部は頂点だけの表現で、相異なる端の
$S,T$ は一回ずつなので、各中間で対応する頂点型を通る成分以外は零である。

両方向がただ一つの未指定 $i$ に集中する場合は、同じ証明で $i$ から一周して $i$ に至る長さ4を取る。
途中三位置のExtは消えており、$a_i=b_i=1$ なので二つの持上げ障害の標的次元は高々2ずつである。
実際この場合に限り頂点との第二Extは $a_i+b_i=2$ である。
経路汎関数空間の次元35より同時に非零解がある。頂点 $i$ は頭と末尾に二度出るが、
正長4の経路は末尾を消し、頭から末尾へ非零に作用する。
従って頭への非分裂矩形と末尾からの非分裂矩形を連接してその間の指定位置を通れる。
$E=0$ から三段を貼れるのでいずれも四行列の恒等式に反し、一向きは全て零である。

もし頂点と直交する非薄単純があれば、その次元は均衡で $c\ge2$、尾部の補題を
いずれかの指定隅に適用して全ての二次元単純も直交すると分かる。なければ
補題\ref{lem:binary-simple-trace-escape}で全単純が次数倍率を除き有限個になり、前節の
GK次元1以下の片側イデアルによる矛盾を得る。これで中心有限性と合成整性が従う。
一般の各隅は同じ次数の可除環であるから実際には全四隅が四行列の恒等式を満たし、
既出の全次数判定が適用される。
\end{proof}
